%% ============================================================
%%  Section: Descriptive Statistics
%%  Depositor Discipline in Modern Banking (2026-04-07 draft)
%%
%%  Changes from 20260406 version:
%%    - Removed all [@@@ ...] placeholders and addressed each one.
%%    - §4.2 (Characteristics) rewritten: removed specific numbers throughout;
%%      retained only the most salient differences (size, deposit structure).
%%      CRE claim corrected: CRE share IS significantly higher for high-CECL
%%      banks in the new table (1.727pp, p<0.05), opposite of what the prior
%%      draft stated.
%%    - §4.2 shortened by roughly one-third. Performance paragraph collapsed
%%      to two sentences; deposit structure paragraph simplified.
%%    - §4.3 (Predictors) shortened to four sentences; removed coefficient
%%      recitation and moved the key takeaway to the opening sentence.
%%
%%  References two tables and two figures defined in tables_figures.tex:
%%    tab:descriptive_sumstats_cecl  (tab01_sumstats_20260403)
%%    tab:cecl_predictors_cross_section  (tab02_cecl_predictors_20260403)
%%    fig:cecl_equity_distribution   (fig01_cecl_dist_20260403)
%%    fig:cecl_adoption_quarter_distribution  (fig02_adoption_dist_20260403)
%%
%%  NOTE: The prior desc-stats file used \label{sec:data} which conflicts
%%  with the data section. This file uses \label{sec:descriptive}.
%% ============================================================

\section{Descriptive Statistics}
\label{sec:descriptive}

In this section, we describe the cross-sectional distribution of the treatment
variable, document the concentration of adoption timing, and compare pre-adoption
characteristics of high- and low-CECL banks.

\subsection{Distribution of the CECL Adjustment and Adoption Timing}
\label{sec:descriptive:distribution}

Figure~\ref{fig:cecl_equity_distribution} plots the cross-sectional distribution
of \texttt{CECL Adj}, the Day-One retained-earnings charge scaled by pre-adoption
book equity, for the community banks in the estimation sample.  The
distribution is right-skewed.  The full-sample mean is 0.28 percent of equity and
the median is essentially zero: roughly half of community banks adopted CECL with
loan portfolios whose forward-looking expected credit losses, under the new
standard, required little incremental reserve above what the incurred-loss
standard had already accumulated.  The right tail is substantial, however.
The threshold for our binary High-CECL indicator, 2.5 percent of equity, sits
at approximately the 95th percentile of the distribution; the 217 banks that
exceed it record a mean adjustment of 4 percent of equity, a material
reduction in reported book capital from a single mandatory accounting event.
Section~\ref{sec:robustness:nonlinearity} examines the full shape of the
dose-response across the distribution, including the left tail of banks with
\emph{negative} (favorable) Day-One adjustments, and shows that the results
are not driven by the single largest adjustment visible in the figure.

Figure~\ref{fig:cecl_adoption_quarter_distribution} shows the distribution of
adoption quarters across the eligible six-quarter window (2022Q3--2023Q4).
Adoption is sharply concentrated in 2023Q1: 3,997 of the 4,320 sample banks
(92.5 percent) appear in their first CECL-compliant Call Report in that quarter,
exactly as prescribed for calendar-year non-public filers.  The remaining
323 banks are distributed across the five adjacent quarters, reflecting
non-calendar fiscal years, early voluntary adoption, and minor filing lags.
This near-uniform adoption calendar strengthens the identification: the
information shock reaches the market in a single concentrated quarter, making
it implausible that the timing itself was chosen in response to bank-specific
deposit conditions.  We nonetheless use bank-specific event time throughout,
assigning each bank its own $t^{*}$, so that the small minority with off-cycle
adoption dates contribute symmetric pre- and post-adoption windows to the panel.

\subsection{Characteristics of High- and Low-CECL Banks}
\label{sec:descriptive:balance}

Table~\ref{tab:descriptive_sumstats_cecl} compares High-CECL and Low-CECL banks
at the adoption cross-section.

High-CECL banks enter the adoption quarter with weaker loan quality and thinner
capital buffers.  Their nonperforming loan ratios and allowance ratios are both
significantly higher, and their CRE concentration is modestly but significantly
greater, consistent with more credit-sensitive loan portfolios requiring larger
expected-loss reserves under the new standard.  Capitalization is meaningfully
lower for high-CECL banks, reflecting the combined effect of elevated credit
exposure and the retained-earnings charge itself.

The pre-adoption deposit structure of the two groups is broadly similar on
dimensions central to our outcome variables.  Uninsured time deposits as a
share of assets do not differ significantly between the groups, and the total
uninsured deposit share is, if anything, slightly \emph{lower} at high-CECL
banks---the opposite of what a run-exposure story would require.  The one
notable asymmetry is that high-CECL banks hold about 1.9 percentage points
more insured time deposits at baseline (9.1 versus 7.3 percent of assets),
reflecting greater reliance on retail certificate funding.  This level
difference does not threaten the design: bank fixed effects absorb it, the
triple-difference compares each bank's uninsured and insured series to each
other within bank and quarter, and the event-study pre-period coefficients on
both deposit series are flat, confirming that the level gap was stable rather
than diverging before adoption.

Performance differences are modest in magnitude and absorbed by bank fixed
effects; they confirm that High-CECL banks represent a coherent group for which
the CECL disclosure carries the most new information about forward-looking credit
risk.

\subsection{What Predicts the CECL Adjustment?}
\label{sec:descriptive:predictors}

Table~\ref{tab:cecl_predictors_cross_section} regresses \texttt{CECL Adj} on
bank characteristics measured in 2022Q4, the last pre-adoption quarter.
Column~(1) uses balance-sheet and credit-quality fundamentals; column~(2)
adds the deposit-composition variables (uninsured, time-deposit, and brokered
shares of deposits) and eight-quarter pre-adoption trends in profitability
and asset quality; column~(3) adds local deposit-market concentration.  The
adjustment loads on credit-risk fundamentals in the expected directions:
larger and less-capitalized banks with higher NPL ratios record larger
charges, and banks on deteriorating profitability trends record larger
charges as well.  Among the funding variables, the time-deposit and brokered
shares enter positively but with economically small coefficients, while the
uninsured deposit share---the variable most relevant to run exposure---is
unrelated to the adjustment.  Unrealized securities losses, the channel
central to the 2023 stress at larger institutions, lose significance once the
expanded controls enter.  The overall fit remains low (adjusted $R^2$ of 1.4
percent in the most saturated specification).  As discussed in
Section~\ref{sec:identification}, we do not treat the low $R^2$ as proof of
exogeneity; rather, column~(2) serves as the first stage for the residualized
treatment used in Section~\ref{sec:robustness}, so that the identifying
variation in that exercise is orthogonal to every predictor in the table by
construction.  Because bank fixed effects in the main regressions absorb all
time-invariant differences, the cross-sectional correlations documented here
do not themselves confound the within-bank treatment estimates.
